% §1 Introduction
% Status: Written (stub — fill in after results)

\section{Introduction}
\label{sec:introduction}

Exascale HPC systems are heterogeneous. Three major GPU vendors --- NVIDIA, AMD, and Intel ---
each expose different memory hierarchies, execution models, and compiler backends.
To write portable code, developers use abstraction layers: Kokkos, RAJA, SYCL, and Julia.
These layers promise portability. They do not always deliver performance.

The current state of practice is guesswork. Developers choose abstractions based on team familiarity,
project history, or vendor recommendation --- not evidence. When performance degrades, root cause
analysis is ad hoc. There is no systematic, measurement-driven framework for making these decisions.

This paper makes the following contributions:

\begin{enumerate}
    \item An \textbf{empirical taxonomy} --- the first systematic classification of abstraction
          failure modes and success patterns across architectures and workloads (Section~\ref{sec:taxonomy}).
    \item A \textbf{measurement methodology} --- a reproducible protocol for cross-layer
          performance attribution (Section~\ref{sec:methodology}).
    \item A \textbf{quantified decision framework} --- evidence-based thresholds and rules
          for abstraction selection (Section~\ref{sec:framework}).
    \item A \textbf{performance portability dataset} --- TODO configurations across 5 abstractions
          $\times$ 3 architectures $\times$ 7 workloads.
    \item A \textbf{decision support tool} (\texttt{abstraction-advisor}) --- a command-line
          advisor that ingests workload profiles and recommends abstraction strategies.
\end{enumerate}

% TODO: Add paper organization paragraph after structure is finalized.
